
\section{Relevo y captura de errores}
Un relevo identifica pasos completados, elementos abiertos, pruebas pendientes, configuración temporal y anomalías. Permite reconstruir el estado sin depender de memoria o conversación. Interrupciones y turnos son condiciones previsibles que los procedimientos y registros deben soportar.

Una segunda firma sin independencia ni evidencia de qué se comprobó puede ser solo apariencia de barrera. El método de captura debe responder al error y a los requisitos de la tarea. La competencia incluye entender el objetivo de la comprobación.
\section{Herramientas y cadena de medida}
Se necesita equipo adecuado al rango y tolerancia, estado de calibración y condiciones de uso. Calibración no declara correcta cualquier lectura: incertidumbre, resolución y técnica también influyen. Si aparece equipo fuera de tolerancia, se evalúan trabajos afectados; cambiar su etiqueta no resuelve medidas anteriores.

El control de herramientas y objetos extraños relaciona inventario, acceso, uso y cierre. Una pérdida puede requerir investigación por el posible efecto sobre producto. Almacenamiento y preservación conservan condiciones de piezas y consumibles, incluida caducidad cuando proceda.
\section{Contener el problema y eliminar su causa}
La acción inmediata contiene; la correctora busca evitar repetición. Retirar una copia incorrecta de manual puede ser necesario e insuficiente: se revisa distribución, selección, alertas y comprobación. La eficacia se demuestra después, no por la fecha de cierre del formulario.

145.A.200 relaciona seguridad y seguimiento de cumplimiento con funciones específicas. Investigar solo quién cometió el último error no explica por qué las barreras permitieron el evento. Esta perspectiva no impone una metodología única; enseña a estudiar condiciones y sistema.
\repaso{¿Un útil calibrado demuestra por sí solo apriete correcto?}{También importan selección, configuración, técnica, dato, secuencia y condiciones. Calibración es una evidencia dentro de una cadena más amplia.}
